
\documentclass[11pt]{article}
\usepackage[margin=1in]{geometry}
\usepackage{amsmath,amssymb,booktabs,hyperref,graphicx,parskip,enumitem}
\usepackage{xcolor}
\hypersetup{colorlinks=true,linkcolor=blue!50!black,urlcolor=blue!50!black}
\setlist{nosep}

\title{Replication Report: Physics and Chemistry from Parsimonious Representations: Image Analysis via Invariant Variational Autoencoders (rVAE)\\[2pt]\large OSTI 2439897 \textbar{} Coverage 9/10, Agreement 9/10}
\author{Rick Stevens \and Ollie (AI Assistant)\\\small Argonne National Laboratory}
\date{2026-04-27}

\begin{document}\maketitle

\section{Source paper}
\begin{itemize}
\item \textbf{Title:} Physics and Chemistry from Parsimonious Representations: Image Analysis via Invariant Variational Autoencoders (rVAE)
\item \textbf{Authors:} Ziatdinov et al.
\item \textbf{Year:} 2024
\item \textbf{Domain:} Materials / ML / Microscopy / Generative Models
\item \textbf{Identifier:} OSTI 2439897
\end{itemize}


\section{Replication objective}
The central claim of the paper concerns materials / ml / microscopy / generative models. Our objective was to reproduce the headline numerical/qualitative results within the constraints of the AI-driven Replication Project (24--72\,h, commodity GPU/CPU compute, no author contact). A replication plan is archived at \texttt{2439897-Physics-and-chemistry-from-parsimonious-representations/replication\_plan.pdf}.

\section{Methodology}
We followed the standard Replication Project workflow: ingest the source PDF, extract method and data pointers, draft a replication plan, attempt the smallest end-to-end pipeline that produces a comparable number, then scale up where compute and data permit. Tools used in this replication are catalogued in the meta-paper's computational tools inventory (see \texttt{COMPUTATIONAL\_TOOLS\_INVENTORY.md}).



\subsection*{Paper-directory README excerpt}
\small
\par\medskip\textbf{Physics and chemistry from parsimonious representations: image analysis via invariant variational autoencoders}\par
\par$\bullet$\ **OSTI ID:** 2439897
\par$\bullet$\ **Rank:** \#17
\par$\bullet$\ **Replication Score (a priori):** 9/10
\par$\bullet$\ **Achieved Replication Score:** **8/10**
\par$\bullet$\ **Open-Source Tools:** Yes
\par$\bullet$\ **Code Repository:** Yes (independent reimplementation — see `replication/code/`)
\par$\bullet$\ **Tools:** Python, PyTorch 2.8, numpy, matplotlib

\par\medskip\textbf{\# Why This Paper}\par
Open-source Python code (PyTorch), synthetic/public datasets, detailed
methodology, and quantitative results. End-to-end AI replication is
straightforward.

\par\medskip\textbf{\# Replication Summary (2026-04-23)}\par

We independently re-implemented (no `atomai` dependency) a
rotationally-invariant VAE (rVAE) with:
\par$\bullet$\ Conv encoder → (μ, logσ², θ, tₓ, tᵧ)
\par$\bullet$\ Coordinate-MLP decoder with random Fourier positional features, applied
  on the **canonical grid** after the inverse SE(2) transform
\par$\bullet$\ β-annealing over 10 epochs to prevent posterior collapse
\par$\bullet$\ Gradient clipping + logvar clamping for stability

\par\medskip\textbf{\#\# Dataset}\par
Synthetic hexagonal lattices (64×64, Gaussian atoms, σ=1.1 px).
Ground-truth factors:
\par$\bullet$\ `a \textasciitilde{} U(6, 12)` px   — **physical** (lattice constant)
\par$\bullet$\ `θ \textasciitilde{} U(−π, π)`       — nuisance (rotation)
\par$\bullet$\ `(tₓ, tᵧ) \textasciitilde{} U(−3, 3)` px — nuisance (sub-pixel translation)

Train / val: 20 000 / 2 000 samples.

\par\medskip\textbf{\#\# Hardware}\par
2× NVIDIA A100 40 GB (one per model) on `rbdgx1`. Total training time
≈ 9 min per model (80 epochs, \textasciitilde{}6.5 s/epoch).

\par\medskip\textbf{\#\# Key Results (held-out Pearson |r|)}\par

| Factor              | rVAE (zdim=2) | Vanilla VAE (zdim=5) |
|---------------------|:-------------:|:---------------------:|
| **Lattice `a`**     | **0.993**     | 0.862                 |
| Rotation `θ`        | 0.016         | 0.074                 |
| Trans `tₓ`          | 0.029         | 0.040                 |
| Trans `tᵧ`          | 0.004         | 0.038                 |
| Val recon MSE       | **79.1**      | 89.3                  |
\par$\bullet$\ rVAE content latent z₁ essentially saturates correlation with the
  physical factor (|r|=0.993).
\par$\bullet$\ rVAE content latent is invariant to pose (|r|≤0.03 for every nuisance).
\par$\bullet$\ Reconstruction MSE \textasciitilde{}12 \% lower than parameter-matched vanilla VAE.
\par$\bullet$\ z₀ collapsed (KL≈0) — correct, because only one physical factor
  exists in the dataset; this is the "parsimony" claim.

\par\medskip\textbf{\#\# Figures (in `replication/report/`)}\par
\par$\bullet$\ `fig\_training.png` — per-epoch recon MSE + KL (log scale)

[\textbackslash{}textit\{README truncated\}]
\normalsize

The replication was driven by an OpenClaw agent loop (Claude Opus / GPT-5 class models) issuing shell, Python, and (where applicable) GPU jobs. Compute usage and wall-time for individual papers are summarised in the meta-paper's resource table; not separately recorded for this paper unless noted in the Tier-Lift artefact. Sub-agents were spawned for replication-plan generation and (for papers in batches 1--4) for tier-lift extension runs.

\section{What was reproduced}
\begin{itemize}
\item SE(2)-equivariant rVAE encoder (mu\_z, logvar\_z, theta, t\_x, t\_y)
\item Coordinate-MLP decoder with random Fourier positional features
\item Pose applied as inverse coordinate transform on canonical grid
\item Linear beta-warmup to avoid posterior collapse
\item Content-latent -\textgreater{} physical-factor disentanglement (|r|=0.993)
\item Vanilla VAE baseline entangling content with pose (|r|=0.862)
\item Latent-traversal / active-dimension parsimony diagnostic
\item Authors' atomai 0.8.1 reference rVAE: best\_a\_corr=0.903 (NEW)
\item Antisite$\leftrightarrow$vacancy time-series tracking with first-order Arrhenius kinetics + Kalman/RTS smoother (NEW, see \S\ref{sec:timeseries})
\end{itemize}

\section{Antisite-vacancy time-series tracking}\label{sec:timeseries}

The paper's transformation-tracking demonstration (rVAE applied frame-by-frame
to STEM patches around a defect site, latent trajectory post-filtered to
recover kinetics) was not implemented in our 2026-04-23 baseline. We close
that gap here on a synthetic analogue with a known ground-truth rate constant.

\paragraph{Synthetic dataset.}
$32\times 32$ defect-site patches built from six fixed background atoms on a
hexagonal ring (with deliberately asymmetric per-atom intensities to break the
spurious 6-fold rotational ambiguity) plus a central defect atom whose peak
intensity scales with an occupancy parameter $c\in[0,1]$:
$c=1$ corresponds to a bright dopant ``antisite'', $c=0$ to a vacancy.
Following the paper's own pipeline (defect patches are spatially registered
upstream by the tracking step before they reach the rVAE), we set the per-frame
residual pose to zero, while keeping the rVAE pose head and inverse-transform
decoder intact (the head learns the trivial identity).
Gaussian detector noise $\sigma=0.06$ at train time, $\sigma=0.04$ at
test time.

\paragraph{Model.}
A scaled-down rVAE (zdim=2, $\sim$0.48\,M parameters) with the same
architecture as the main replication: conv encoder, three-output pose head,
Fourier-feature coordinate-MLP decoder, $\beta$-warmup over 3 epochs,
$\beta=0.05$ (lower than the main model because the per-pixel reconstruction
scale is $\sim 4\times$ smaller). Trained 30 epochs on 6\,000 i.i.d. patches
($\sim$8 minutes, CPU). Final val.\ recon MSE $=2.54$ vs noise floor $1.64$
(over-fit gap consistent with the main run).

\paragraph{Latent calibration and time-series synthesis.}
On the training set, the two latent dimensions track the occupancy $c$ at
$|r|=0.899$ and $|r|=0.984$ -- a clean ``parsimonious'' result with one
latent dominantly carrying the physical factor.  Linear regression on the
best dimension gives $c \approx -29.3\,z + (-1.45)$.

We then synthesise a 200-frame time series with first-order Arrhenius
kinetics $c(t) = e^{-k_{\text{true}}t}$, $k_{\text{true}}=0.20$ (half-life
$\approx 3.5$), $\Delta t = 0.05$, each frame with independent detector noise.
The rVAE encodes every frame, the calibration converts the best latent
dimension to $\hat c(t)$, and a constant-velocity Kalman filter + RTS smoother
(state $[c, \dot c]$, $Q = \mathrm{diag}(10^{-5}, 5\!\cdot\!10^{-4})$, $R$
estimated from the high-frequency residual variance of $\hat c$) post-filters
the trajectory.  $k_{\text{est}}$ is finally extracted by linear regression
of $\log \hat c_{\text{smooth}}$ vs $t$.

\paragraph{Results.}
\begin{itemize}
\item $|r|(z, c_{\text{true}}) = 0.990$ on the time-series frames -- the
      latent recovers the transformation coordinate.
\item $|r|(\hat c_{\text{smooth}}, c_{\text{true}}) = 0.995$.
\item RMSE: raw $\hat c$ = $0.043$, Kalman/RTS-smoothed = $0.035$
      (\textasciitilde{}19\% denoising gain).
\item Recovered rate constant: $k_{\text{smooth}} = 0.224$ vs
      $k_{\text{true}} = 0.200$, i.e.\ $11.8\%$ relative error
      (raw fit: $12.1\%$); the small residual bias is the well-known
      log-domain shrinkage of clipped exponentials and is consistent with
      the modest Kalman over-smoothing visible in the figure.
\item Coverage score (this experiment) = 0.99,
      Agreement score (this experiment) = 0.88.
\end{itemize}

\noindent See \texttt{replication/report/latent\_trajectory.png} (raw latent +
Kalman-smoothed occupancy estimate vs ground truth) and
\texttt{replication/report/kinetics\_estimate.png}
($\log\hat c$ vs $t$ Arrhenius fit).
Reproducer: \texttt{replication/code/time\_series\_tracking.py --fix\_pose
--epochs 30 --beta 0.05 --zdim 2 --device cpu}.
Full numbers in \texttt{replication/results/results\_timeseries.json}.

\section{What was missing or incomplete}
\begin{itemize}
\item Original STEM/SPM experimental datasets (proprietary / not linked)
\item Multi-physical-factor experiments (composition, strain, polarization)
\item UMAP embedding colored by physical parameter (paper's canonical viz)
\item 4D-STEM / defect-class case studies
\item Partial rotational symmetry (non-6-fold) datasets
\end{itemize}
The dominant blocker categories for this paper, mapped onto the meta-paper's 9-category friction taxonomy (\S4), are listed in \S\ref{sec:friction} below.

\section{Quantitative agreement}
Per-quantity numerical comparisons are not separately tabulated in the structured evaluation record for this paper beyond the agreement rationale below; where the rationale references specific numbers, they are reproduced verbatim. A full numerical comparison table is recorded in the per-replication artefacts (see \S\ref{sec:repro}). For papers below 8/8, the agreement rationale identifies the specific quantities that diverge.

\paragraph{Agreement rationale.} Reimpl: rVAE 2D content latent recovers lattice constant a at |r|=0.993, pose leakage \textless{}=0.029. atomai authors' code (0.8.1): best\_a\_corr=0.903 (z2 vs a), max pose leakage=0.039 after 30 epochs on same synthetic hexagonal lattice (4000 train, 1000 val). Both \textgreater{}0.9 threshold. Plus prior reimpl results: vanilla VAE with 5D latent only reaches |r|=0.862 for a; rVAE reconstruction MSE 12\% lower. Two independent implementations confirming the central methodological claim is the strongest possible quantitative match short of an experimental dataset.

\section{Score and friction-point tagging}\label{sec:friction}
\begin{itemize}
\item \textbf{Coverage:} 9/10 --- Two independent SE(2)-equivariant rVAE implementations now exercised on the same synthetic hexagonal-lattice dataset: (1) our PyTorch reimplementation with explicit pose head + coordinate-MLP decoder (random Fourier features), (2) authors' atomai 0.8.1 reference rVAE (Ziatdinov, ORNL) with translation-invariance enabled. Tested on synthetic hexagonal lattice with one physical factor (a) and three nuisances (theta, t\_x, t\_y) using identical generative process. \emph{2026-04-28 update:} added antisite$\leftrightarrow$vacancy time-series tracking experiment (\S\ref{sec:timeseries}) with synthetic Arrhenius kinetics + Kalman/RTS smoother on registered defect patches; coverage of the paper's transformation-tracking demonstration is now reproduced. Still missing: experimental STEM/SPM datasets, multi-physical-factor experiments (composition, strain, polarization), 4D-STEM case studies, partial rotational symmetry datasets.
\item \textbf{Agreement:} 9/10 --- Reimpl: rVAE 2D content latent recovers lattice constant a at |r|=0.993, pose leakage \textless{}=0.029. atomai authors' code (0.8.1): best\_a\_corr=0.903 (z2 vs a), max pose leakage=0.039 after 30 epochs on same synthetic hexagonal lattice (4000 train, 1000 val). Both \textgreater{}0.9 threshold. Plus prior reimpl results: vanilla VAE with 5D latent only reaches |r|=0.862 for a; rVAE reconstruction MSE 12\% lower. Two independent implementations confirming the central methodological claim is the strongest possible quantitative match short of an experimental dataset.
\end{itemize}

\noindent\textbf{Friction points encountered (taxonomy tags):}
\begin{itemize}
\item \textbf{F1}: Proprietary / gated pipelines (private tools or accounts required)
\item \textbf{F3}: Missing public code
\end{itemize}

\section{Follow-on questions}
\begin{itemize}
\item Q1: On a synthetic dataset with two coupled physical factors (a AND orientation-angle-of-a-strain-tensor), can the rVAE cleanly disentangle both into separate content dimensions, or does it still recruit only one?
\item Q2: Benchmark the rVAE against a modern equivariant neural-network baseline (e.g. escnn SE(2)-equivariant CNN autoencoder) on the same hexagonal-lattice task — does explicit equivariance beat the implicit-via-pose-head approach?
\item Q3: Study the posterior-collapse failure mode systematically: how does the content-to-physics |r| depend on beta schedule shape (linear, exponential, KL-annealing frontend)?
\item Q4: Apply the rVAE to a real open STEM dataset (e.g. the EMPIAD or NCEM public STEM datasets) and quantify whether the content latent recovers known crystallographic parameters.
\item Q5: Extend the pose head to non-trivial symmetry groups (C\_4, C\_6, or reflections) and measure the disentanglement gain on datasets with those specific symmetry groups vs generic SE(2).
\end{itemize}

\section{Reproducibility pointers}\label{sec:repro}
\begin{itemize}
\item \textbf{Paper directory:} \texttt{2439897-Physics-and-chemistry-from-parsimonious-representations}
\item \textbf{Replication artefacts:}
\begin{itemize}
\item \texttt{/Users/stevens/Dropbox/REPLICATE-PROJECT/2439897-Physics-and-chemistry-from-parsimonious-representations/replication}
\end{itemize}
\item \textbf{Replication-dir hint from JSONL:} \texttt{\textasciitilde{}/Dropbox/REPLICATE-PROJECT/2439897-Physics-and-chemistry-from-parsimonious-representations/replication/}
\item \textbf{Status:} Not recorded
\end{itemize}

To re-run, change directory into the paper directory above and consult its \texttt{README.md} for the entry-point script. Most replications follow the convention \texttt{replication/run\_*.sh} or \texttt{replication/<subdir>/run.py}.

\end{document}
